\documentclass[12pt]{article}
\usepackage[margin=1in]{geometry}
\usepackage{enumitem}
\usepackage{titlesec}
\usepackage{parskip}
\usepackage{graphicx}
\usepackage{hyperref}
\usepackage{fontenc}

\title{\textbf{Lipset and Rokkan 1967 RG Notes}\\
\large Lipset, Seymour Martin, and Stein Rokkan. ``Cleavage Structures, Party Systems, and Voter Alignments: An Introduction.''\\
In \textit{Party Systems and Voter Alignments: Cross-National Perspectives}, eds. Seymour Martin Lipset and Stein Rokkan.\\
New York: Free Press, 1967.}
\author{}
\date{}

\begin{document}

\maketitle

\section{Main Idea}



%--------------------------------------------------
\section{General Notes}
%--------------------------------------------------

\begin{itemize}[leftmargin=*, itemsep=4pt]
    \item Contemporary party divides are primarily the result of the consequences from two major historical transitions: the national revolution, and the industrial revolution. Party populations have been largely frozen around their primary-served groups following these transitions, demonstrating a great deal of stability with respect to their primary constitutent bases. This is largely determined by the order and combination of the above junctures and their effects, rather than by contemporary considerations.
    \item Four cleavages generated by these revolutions, upon which party ``identities'' have largely been based:
    \begin{enumerate}
        \item Center vs. periphery
        \item State vs. church
        \item Land vs. industry
        \item Owner vs. worker
    \end{enumerate}
    \item Cleavages require three components in order to become politically consequential or relevant to political identity:
    \begin{enumerate}
        \item Empirical basis: there must be some sort of identifiable social division, a way to distinguish the cleavage itself and include/exclude individuals based on their characteristics
        \item Normative element: there must be some set of shared identity, ideals, or interests, such that those values can legitimate the division between the ``out''-groups
        \item Organizational expression: groups, parties, institutions, etc. which mobilize and represent the division (this is the difference between cleavages and plain social difference)
    \end{enumerate}
    \item A political movement must cross four thresholds to warrant full institutionalization as a party:
    \begin{enumerate}
        \item Legitimation: members can petition, criticize, and organize as an opposition group within the broader political Systems
        \item Incorporation: legislative suffrage is extended to the movement's supporters
        \item Representation: the movement can, in fact, win seats
        \item Majority power: the movement can plausibly compete for a position in the government, and is not a perpetual minority voice
    \end{enumerate}
    \item Cleavages can either be cross-cutting or reinforcing:
    \begin{itemize}
        \item Cross-cutting cleavages occur when some class or group defined by a cleavage is itself divided by another cleavage, resulting in greater party fragmentation and cross-pressurization of voters
        \item Reinforcing cleavages occur where cleavages are largely aligned, such that subgroups are not further split and group identity is reinforced. This tends to create polarization along a smaller number of much deeper lines
    \end{itemize}
\end{itemize}


\section{Important Specific Factual Claims}

\begin{itemize}[leftmargin=*, itemsep=4pt]
    \item Use of an A/G/L/I spatial model, but somewhat ill-defined and difficult to work through...
    \item Parties have an incentive/mechanism for their own survival: once established, they will evolve and adapt such that some version of the institution itself can survive. This becomes a foundation of path dependency; it is difficult to change directions once the party has been established
\end{itemize}


\section{Connections}

\begin{itemize}
    \item General connection on group identity -- see Anderson 1983
    \item Many of the legitimation mechanisms seem to tie back to Powell 2000 and a question of majoritarian governance -- a proportional approach may allow for a group to become politically relevant with much greater accessibility
    \item Another connection to Dahl: is the political system inclusive, and is the opposition meaningfully part of the political process
    \item Similar arguments to Moore 1966's characterization of political systems as a reflection of historical orders of wealth and welfare, though along generally different approaches
\end{itemize}


\section{My Critiques}

\begin{itemize}
    \item Is this still true? Factually, it seems like it may have fallen apart quite a bit, particularly as left-of-center parties have aligned around more educated and wealthier voters. It is also a tad difficult to truly consider existing party conditions in 1920s Europe given the presence (and fate) of the Weimar Republic
\end{itemize}


\end{document}
